\documentclass{article}
\usepackage[T1]{fontenc}
\usepackage{lmodern}
\usepackage{graphicx}
\usepackage{amsmath,amssymb}
\usepackage[a4paper,margin=2cm]{geometry}
\usepackage[absolute,overlay]{textpos}
\setlength{\TPHorizModule}{1mm}
\setlength{\TPVertModule}{1mm}
\usepackage[indonesian]{babel}
\usepackage{hyperref}
\setlength{\emergencystretch}{2em}
% unit: Halaman 22

\begin{document}

% === Halaman 22 (python/native) ===

• OTP Tidak cocok untuk komunikasi dinamis. Bayangkan untuk komunikasi 

modern seperti:

o WhatsApp 

o HTTPS 

o VPN 

o cloud storage 

o email 

o komunikasi IoT, dll 

   data yang dikirim terus berubah, perangkat bisa berganti, pengguna bisa 

offline, dan pesan memiliki ukuran yang berbeda-beda.

• Dengan OTP, setiap tambahan data membutuhkan tambahan kunci yang 

benar-benar acak dan belum pernah digunakan. Ini membuat manajemen 

kunci menjadi sangat berat.

22

\end{document}
